\documentclass[11pt]{article}

\usepackage[T1]{fontenc}
\usepackage[utf8]{inputenc}
\usepackage{lmodern}
\usepackage[margin=1in]{geometry}
\usepackage{amsmath,amssymb,amsthm,mathtools}
\usepackage{booktabs,array}
\usepackage[shortlabels]{enumitem}
\usepackage{xcolor}
\usepackage{hyperref}
\usepackage{xurl}
\usepackage{microtype}
\usepackage[numbers,sort&compress]{natbib}
\newcommand{\OPHPaperReleaseID}{r2011}
\newcommand{\OPHPaperReleaseDate}{August 1, 2026}
\newcommand{\OPHPaperReleaseBanner}{%
\begin{center}
\small\textbf{Paper release:} \texttt{\OPHPaperReleaseID}\quad\textbf{Released:} \OPHPaperReleaseDate
\end{center}
}
\newcommand{\PaperReleaseID}{\OPHPaperReleaseID}
\newcommand{\PaperReleaseDate}{\OPHPaperReleaseDate}
\newcommand{\PaperReleaseBanner}{%
\begin{center}
\small\textbf{Paper release:} \texttt{\PaperReleaseID}\quad\textbf{Released:} \PaperReleaseDate
\end{center}
}


\hypersetup{colorlinks=true,linkcolor=blue,citecolor=blue,urlcolor=blue,
  pdftitle={Verified Projection-Event Calculus in Lean 4},
  pdfauthor={Bernhard Mueller}}
\setlength{\parindent}{0pt}
\setlength{\parskip}{0.55em}
\setlength{\emergencystretch}{3em}
\renewcommand{\arraystretch}{1.12}
\newcolumntype{L}[1]{>{\raggedright\arraybackslash}p{#1}}
\renewcommand{\_}{\textunderscore\allowbreak}

\newtheorem{theorem}{Theorem}
\newtheorem{proposition}{Proposition}
\newtheorem{definition}{Definition}
\theoremstyle{remark}
\newtheorem{remark}{Remark}

\newcommand{\C}{\mathbb C}
\newcommand{\Mn}{\mathrm{M}_n(\C)}
\newcommand{\Tr}{\operatorname{Tr}}
\newcommand{\one}{\mathbf 1}
\newcommand{\pinch}{\mathcal E}
\newcommand{\avg}{\mathcal A}
\newcommand{\lued}{\mathcal L}
\newcommand{\norm}[1]{\lVert #1\rVert}
\title{\textbf{Verified Projection-Event Calculus in Lean~4:\\
Bundled Arbitrary-Partition Pinching, L\"uders Retractions,\\
and CHSH Interoperability}}
\author{Bernhard Mueller\thanks{Corresponding author:
\href{mailto:bernhard@floatingpragma.ai}{\texttt{bernhard@floatingpragma.ai}}}}
\date{\OPHPaperReleaseDate}

\begin{document}
\maketitle
\OPHPaperReleaseBanner

\begin{center}
\small\textbf{Author affiliation:} Bernhard Mueller, Pragma Research Inc.
\end{center}

\begin{abstract}
Finite quantum measurement looks simple on paper: multiply matrices, take a
trace, and normalize. In Lean, each step crosses a different interface and
every side condition must be supplied. This paper presents a Lean~4 library for
projection events over finite complex matrix algebras. A projective partition
defines a star-subalgebra of matrices that commute with its projectors and a
complex-linear pinching map onto that subalgebra. The map is unital, positive,
trace preserving, and idempotent; its range and fixed points are exactly the
partition commutant. It also satisfies the bimodule law,
Hilbert--Schmidt Pythagorean and contraction identities, and a
trace-duality uniqueness theorem.
It is also exactly the positive uniform random-unitary average over all
independently signed block reflections. A global-sign control shows why
independent signs are necessary.

A second expectation averages onto the commutative span of the partition.
The two maps satisfy tower laws, preserve the partition's Born statistics,
and commute with conditioning on a partition member. L\"uders conditioning
is available as a total matrix formula and a typed state update, with an
exact fixed-state characterization. An operator-norm estimate lifts the
Tsirelson theorem to a state-level
Clauser--Horne--Shimony--Holt bound for projection events, and a declared
maximally entangled witness on a checked slot-locality interface attains the
    bound $2\sqrt2$ exactly, while every fixed record-diagonal state read
    through four observables diagonal in the same basis and valued in
    $[-1,1]$ obeys the classical bound $2$. An ideal static phase-POVM/count
    model preserves the exact tomography and blindness delimitations, and an
    explicit synthetic inhabitant proves that exact fit alone is not an
    operational instrument and does not discharge the phase-sensitive-effect
    premise. The artifact pins
its Lean and Mathlib revisions and records declaration-level axiom
dependencies. Complete positivity, a Physlib correspondence, and rank-one
classicality lie outside the formalization. A
separate support countermodel proves that totalizing the matrix logarithm at
zero cannot define support-aware relative entropy: the raw trace formula
vanishes on an orthogonal pure-state pair.
\end{abstract}

\noindent\textbf{Keywords:} formal verification, Lean~4, projective measurement,
partition pinching, support-aware relative entropy, L\"uders update, CHSH

\medskip
\noindent\textbf{Mathematics Subject Classification (2020):}
Primary 68V20; Secondary 68V15, 81P15, 81P40, 46L53.

\section{Introduction}
\label{sec:intro}

Finite quantum measurement is short on paper. A proof calls a matrix a state,
inserts a projection, normalizes the result, and moves on. Lean asks for every
condition skipped in that sentence. The matrix must be positive and have unit
trace. The projection must be Hermitian and idempotent. Normalization needs a
nonzero denominator. Later norm estimates require a further set of
instances.

Machine-checked quantum information is an active area, spread over several
systems and several design philosophies. Mathlib \cite{mathlib2020}, the
central mathematics library of Lean~4 \cite{lean4}, contains an order-form
CHSH/Tsirelson theorem and the \texttt{IsCHSHTuple} interface
\cite{mathlibCHSH2026}. Physlib's \texttt{QuantumInfo.Finite.Pinching}
constructs the spectral pinching of a state as a completely positive
trace-preserving map and proves commutation, fixed-point, Pythagorean, and
pinching-bound results \cite{physlib,physlibPinching2026}. Lean-Quantum
provides basis-independent infrastructure for states, channels, tensor
products, Choi and Kraus representations, and trace inequalities
\cite{kasauraEtAl2026LeanQuantum}; Lean-QIT develops a compositional
infrastructure for quantum-information theory \cite{zhuEtAl2026LeanQIT}. Zhao
and Yu formalize CHSH rigidity in Lean~4 \cite{zhaoYu2026CHSHRigidity}.
Isabelle/HOL developments cover projective measurement, quantum computation,
and the CHSH upper bound
\cite{echenim2021afp,echenimmhalla2023afp,echenimmhalla2024,bordg2021}.
These developments choose different carriers for states and different
bundling conventions, and results proved in one seldom apply verbatim in
another. The engineering question underneath is interface design: which
objects stay bare predicates, which get bundled, and how finished results are
handed to the interfaces a large library expects. We make no priority claim
for projective measurement, pinching, or Tsirelson's inequality.

The library described here packages the side conditions of finite projective
measurement into a projection-event API over Mathlib matrices. Its main input
is an arbitrary projective partition of the identity. The partition induces
two canonical expectations: the pinching, which deletes off-block terms and
maps the matrix algebra onto the partition commutant, and the average, which
projects further onto the commutative span of the partition. The same event
definitions feed the L\"uders update and the CHSH adapter. This keeps
algebraic identities separate from results that use trace, positivity, or
state normalization.

The mathematics is classical: L\"uders introduced the selective update rule
\cite{lueders1950,buschlahti2009}; pinching by a family of orthogonal
projections is a standard averaging operation in matrix analysis
\cite{bhatia2000} and in quantum information \cite{hayashi2017}; the bimodule
and trace-duality laws proved here are the finite-matrix forms of the
defining properties of conditional expectations onto subalgebras
\cite{umegaki1954,takesaki1972}; and the CHSH combination with its quantum
bound goes back to Clauser, Horne, Shimony, and Holt \cite{chsh1969} and to
Tsirelson \cite{tsirelson1980}. The contribution of the library is the
checked interface: definitions bundled the way Mathlib expects, exact range
and uniqueness theorems for both expectations, and adapters into existing
Mathlib results.

Lean declaration identifiers are printed throughout in monospaced type, for
example \texttt{partitionPinching\_idem}; the underscore is part of the
identifier. To locate a declaration, search for the exact printed identifier,
after removing only TeX's protective backslash before each underscore, in the
\texttt{EventAlgebra/*.lean} files. In the public repository these files are
under \texttt{Lean/EventAlgebra/}. Table~\ref{tab:modules} maps each module
name to its subject matter. The top-level \texttt{EventAlgebra.lean} file
imports these fifteen modules alongside the wider OPH event-algebra library.

The checked development contributes fifteen pieces:

\begin{enumerate}[(i)]
\item an arbitrary supplied projective partition, independent of a selected
state or spectral resolution, with its commutant bundled as a
\texttt{StarSubalgebra} and its span bundled with multiplication, star,
commutativity, and centrality theorems;
\item partition pinching bundled as a complex \texttt{LinearMap}, with exact
range and fixed-point theorems, positivity, unitality, trace preservation, the
commutant bimodule law, Hilbert--Schmidt geometry, and map-level uniqueness;
\item an exact random-unitary representation of arbitrary partition pinching
as the positive uniform average over all independently signed block
reflections, together with a global-sign negative control;
\item an orthogonal-state support countermodel showing that the totalized
matrix logarithm returns zero where any support-aware extended relative
entropy must return infinity;
\item partition averaging bundled as a complex \texttt{LinearMap} onto the
span, with exact range, trace duality, map-level uniqueness, tower laws
against the pinching, preservation of partition Born statistics, and a
classical-conditioning collapse theorem;
\item a typed nonzero-weight L\"uders state update, a support theorem and an
unguarded fixed-point characterization that exhibit conditioning as a
retraction onto its certainty set, and naturality of the typed update under
partition pinching for commutant events;
\item an event-to-observable CHSH client that discharges Mathlib-compatible
selfadjointness, involutivity, and cross-commutation hypotheses before
applying a complementary norm-form bound;
\item a checked state-expectation bound that converts the norm-form theorem
into a state-level CHSH bound;
\item a finite-state Robertson theorem for supplied density matrices and
Hermitian observables, with the ordinary-commutator form, a noncommuting
saturation control, and a separate zero-variance control;
\item the exact quotient induced by a supplied projective partition: equality
after block pinching is equivalent to equality under every commutant trace
test, while block pinching and commutative partition averaging both erase
cross-sector corners;
\item an exact Tsirelson saturation witness: a declared maximally entangled
state and four declared projection events inhabiting a bipartite
slot-locality interface of the wider library, whose CHSH expectation equals
$2\sqrt2$ exactly, packaged with the cited upper bounds, together with a
record-diagonal delimitation proving the classical bound $2$ for every
diagonal state read against four diagonal observables valued in $[-1,1]$;
\item an ideal static phase-POVM/count-fit structure with conditional
tomography, exact diagonal/real-closure blindness, and an explicit synthetic
inhabitant proving that the structure does not type source operations,
producer receipts, or discharge of the phase-sensitive-effect premise;
\item an exact determination boundary for that structure: a coordinate
characterization of the phase-sensitivity clause, run-literal pinning of the
state diagonal with one free off-diagonal coordinate, affine dependence of
every count frequency on that coordinate, a decidable integer receipt window
on the phase counts, and the minimal phase-mass receipt
(Section~\ref{sec:phase-determination});
\item a static conformance layer for that structure: taking the committed exact
lift as a declared phase-sensitive effect, a
deterministic exact-arithmetic calculator, an independent verifier, and a
composed Lean arithmetic receipt inhabit the static model with generated
expected-frequency numerators satisfying the assumed Born-fit equations,
the scaling law, and the integer receipt window
(Section~\ref{sec:phase-attainment}); and
\item a reproducible, version-pinned artifact with per-declaration axiom
output, no proof placeholders, and a documented catalog of the Mathlib gaps,
scoped-instance hazards, and workarounds met during construction
(Section~\ref{sec:mathlib}).
\end{enumerate}

Two results are outside the current library. No theorem identifies an
arbitrary partition with Physlib's spectral partition. There is no rank-one
classical-representation theorem. The CHSH bounds carry a saturation witness
(Section~\ref{subsec:saturation}); the witness state and events are declared
data, produced by no theorem from any source object. A companion module,
\texttt{Dynamics/ChoiCPTP.lean}, defines an explicit finite-matrix CPTP
predicate and proves both partition expectations completely positive and
trace preserving; the fifteen-module audit here covers their event-algebra
interfaces rather than recounting that companion development.

Section~\ref{sec:setting} fixes the mathematical and Lean interfaces.
Sections~\ref{sec:events} and \ref{sec:lueders} treat projection events and
selective update. Section~\ref{sec:pinching} gives the bundled partition
pinching and its geometry, Section~\ref{sec:spectral-boundary} the
random-unitary representation and logarithm-support boundary, and
Section~\ref{sec:average} the averaging
expectation, the two-level structure, and the partition-relative operational
quotient. Sections~\ref{sec:functional}, \ref{sec:robertson}, and
\ref{sec:chsh} describe the state functional, finite-state uncertainty, and
the CHSH clients with the exact saturation witness and its record-diagonal
delimitation. Section~\ref{sec:phase-instrument} describes the ideal static
phase-POVM/count-fit model, its synthetic non-discharge control, and the
declared-effect conformance fixture for the phase-count layer.
Section~\ref{sec:mathlib} records the Mathlib engagement.
Section~\ref{sec:evaluation} reports the audit methodology, and
Section~\ref{sec:comparison} positions the development feature by feature.

\section{Mathematical and formal setting}
\label{sec:setting}

Let $\Mn$ be the complex $n\times n$ matrices, with conjugate transpose
$X\mapsto X^*$, trace $\Tr$, and identity $\one$. Following the standard
finite-dimensional setting \cite{vonneumann1932}, a projection event is a
Hermitian idempotent $P=P^*=P^2$, a state is a positive-semidefinite matrix
$\rho$ with $\Tr(\rho)=1$, and the trace pairing is
$\mu_\rho(M)=\Tr(\rho M)$.

Lean represents these interfaces as predicates over Mathlib matrices:
\texttt{IsEvent P} is the conjunction $P^*=P$ and $P^2=P$;
\texttt{IsState rho} is positive semidefiniteness and unit trace. For
functions that return states, the library defines the subtypes
\texttt{ProjectionEvent n} and \texttt{StateMatrix n}. The split is
practical. Algebraic lemmas use raw matrices. A function advertised as a state
update returns a value whose type carries the state conditions.

Table~\ref{tab:modules} lists the fifteen source modules. The comparison in
this work concerns their compiled interfaces and uses, not the number of
declarations in each file.

\begin{table}[ht]
\centering
\small
\begin{tabular}{@{}L{3.0cm}L{9.2cm}@{}}
\toprule
Module & Principal interface \\
\midrule
\texttt{Basic} & projection events, states, trace pairing, closure and Born
weight bounds \\
\texttt{Lueders} & totalized matrix formula, typed state update, repeatability,
composition, support, and fixed points \\
\texttt{Partition\allowbreak Pinching} & projective partitions, bundled commutant,
bundled linear pinching, exact range, geometry, bimodule law, and naturality \\
\texttt{PartitionAverage} & partition span with closure, commutativity, and
centrality; bundled averaging expectation with exact range, uniqueness, tower
laws, statistics preservation, and classical conditioning \\
\texttt{StateExpectation} & trace pairing bundled as a complex-linear
functional, normalization and positivity \\
\texttt{Robertson} & supplied-state Robertson inequality, ordinary-commutator
identity, variance nonnegativity, and exact Pauli controls \\
\texttt{Superselection} & complete operational quotient for partition
pinching and cross-sector invisibility for both declared readouts \\
\texttt{Record\allowbreak Majorization} & positive uniform independent-sign
random-unitary representation of arbitrary partition pinching and a
global-sign negative control \\
\texttt{Spectral\allowbreak Entropy\allowbreak Boundary} & exact support-failure density witness and
totalized-matrix-logarithm no-go for support-aware relative entropy \\
\texttt{Tsirelson} & CHSH ring identity, norm theorem, Mathlib interoperability,
and event-level matrix client \\
\texttt{ExpectationBound} & state-expectation operator-norm bound and the
state-level CHSH corollary \\
\texttt{Tsirelson\allowbreak Saturation} & declared Bell witness on the
slot-locality interface attaining $2\sqrt2$ exactly, packaged saturation
receipt, and record-diagonal classical delimitation \\
\texttt{Operational\allowbreak Phase\allowbreak Instrument} & ideal static
phase-POVM/count-fit structure, conditional tomography, exact blindness
receipts, and a synthetic non-discharge counterreceipt \\
\texttt{Phase\allowbreak Instrument\allowbreak Determination} & coordinate
characterization of the phase clause, run-literal state pinning, affine
count-frequency dependence, decidable integer receipt window, and the
minimal phase-mass receipt \\
\texttt{Operational\allowbreak Phase\allowbreak Attainment} &
declared-effect inhabitant of the count-fit structure with generated
expected-frequency numerators, stored Born-fit equations, scaling consequences,
and the receipt-window clause \\
\bottomrule
\end{tabular}
\caption{Module structure of the checked development.}
\label{tab:modules}
\end{table}

Three Mathlib scopes affect compilation. \texttt{ComplexOrder} equips
$\C$ with the partial order in which $0\le z$ means that $z$ is real and
nonnegative. \texttt{MatrixOrder} supplies the Loewner order.
\texttt{Matrix.Norms.L2Operator} activates the finite-matrix operator norm and
C*-ring instances used by the matrix CHSH theorems. Without the right scope,
Lean may report an instance problem even though the required theorem is in
scope; Section~\ref{sec:mathlib} describes the failure modes.

Documentation tags classify results as \emph{algebra-only} or
\emph{trace-dependent}. These tags do not define an import hierarchy. They
record whether a proof uses ring, star, and norm structure alone, or also uses
$\Tr(\rho M)$, positivity, or normalization.

\section{Projection events and the trace pairing}
\label{sec:events}

The event layer proves the expected closure rules. Zero, identity, and
$\one-P$ are events. Orthogonality is symmetric. Orthogonal sums and products
of commuting events are events. Subevent absorption holds on both sides, and
every event is positive semidefinite. The declaration
\texttt{IsEvent.one\_sub\_two\_smul\_involution} proves that
$\one-2P$ is a selfadjoint involution. The CHSH adapter in
Section~\ref{sec:chsh} uses this result.

The trace pairing is additive, homogeneous, compatible with finite sums, and
obeys the sandwich identity
\[
  \Tr(P\rho P)=\Tr(\rho P)
\]
for idempotent $P$. Hermiticity of $\rho$ and $P$ implies that the pairing is
real. Positivity of $\rho$ and the event property imply
$0\le\mu_\rho(P)$ in Mathlib's order on $\C$; for a state the upper bound is
$\mu_\rho(P)\le 1$. The order-form statements are strictly stronger than
their real-part corollaries, and both forms are provided.

The declaration
\texttt{isEvent\_add\_and\_bornWeight\_add\_of\_orthogonal} returns two facts:
the orthogonal sum is an event, and its Born weight is the sum of the two Born
weights. The event and orthogonality hypotheses carry the first conjunct; the
weight identity is linear and uses none of them.

\section{L\"uders update as a partial state interface}
\label{sec:lueders}

For raw matrices, the library uses the totalized formula
\[
 \lued_P(\rho)=\mu_\rho(P)^{-1}P\rho P.
\]
Lean's field inverse maps zero to zero, so the formula is total and covers
algebraic identities in the zero-weight case. Its return type is a raw matrix
and carries no state guarantee. The typed operation is
\[
\begin{aligned}
\texttt{luedersStateUpdate}:\quad&
 \texttt{StateMatrix n}\to\texttt{ProjectionEvent n}\\
 &\to(\mu_\rho(P)\ne0)\to\texttt{StateMatrix n}.
\end{aligned}
\]
This operation assumes \texttt{NeZero n}. The unique $0\times0$ matrix cannot
have unit trace, so a state-returning function must exclude that dimension.

\begin{proposition}[Checked L\"uders interface]
\label{prop:lueders}
For a state $\rho$, event $P$, and nonzero outcome weight:
\begin{enumerate}[(i)]
\item the raw formula is positive semidefinite and has trace one
(\texttt{luedersUpdate\_isState});
\item the updated state assigns weight one to $P$
(\texttt{bornWeight\_luedersUpdate\_self});
\item update is an idempotent retraction
(\texttt{luedersUpdate\_idem});
\item updates by commuting events compose and commute
(\texttt{luedersUpdate\_luedersUpdate\_of\_commute} and
\texttt{luedersUpdate\_comm});
\item when the state commutes with the event, the update is the normalized
block restriction $\mu_\rho(P)^{-1}\rho P$
(\texttt{luedersUpdate\_of\_commute}); and
\item the fixed-state equivalence
\[
  \lued_P(\rho)=\rho \quad\Longleftrightarrow\quad \mu_\rho(P)=1
\]
holds without a separate nonzero-weight hypothesis
(\texttt{luedersUpdate\_eq\_self\_iff}).
\end{enumerate}
\end{proposition}

The fixed-point theorem needs no nonzero-weight guard. A state fixed by a
zero-weight totalized update would equal the zero matrix, which has the wrong
trace. The raw function covers the zero-weight case; the state theorem
does not carry a redundant premise.

The module assembles these results into a fixed-point description of
conditioning. The set \texttt{certainStates} collects the states that assign
weight one to $P$. Conditioning any state of nonzero weight lands in that set
in a single step (\texttt{luedersUpdate\_mem\_certainStates}), acts on the
set as the identity, and, among states, fixes exactly its elements.
Conditioning on $P$ is therefore a retraction of the nonzero-weight states
onto its own fixed-point set. The identity case rests on a support theorem,
\texttt{mul\_eq\_self\_of\_bornWeight\_one}: a state certain of $P$ absorbs
the event on both sides, $\sigma P=P\sigma=\sigma$. The checked proof
compresses $\sigma$ by the complement event, whose weight vanishes; positive
semidefiniteness forces the zero-trace compression to vanish; and the
Cauchy--Schwarz property of the state's quadratic form eliminates
$\sigma(\one-P)$ itself.

\section{Bundled arbitrary-partition pinching}
\label{sec:pinching}

A \texttt{ProjectivePartition n k} consists of projections
$(P_i)_{i<k}$, pairwise orthogonality, and $\sum_iP_i=\one$. It induces
\[
 \pinch_\pi(X)=\sum_iP_iXP_i,
 \qquad
 N_\pi=\{X:XP_i=P_iX\text{ for every }i\}.
\]
The code defines $N_\pi$ as
\texttt{ProjectivePartition.commutant}, a
\texttt{StarSubalgebra} defined using Mathlib's centralizer, and bundles
$\pinch_\pi$ as \texttt{partitionPinchingLinearMap}. The membership theorem
\texttt{mem\_commutant\_iff} reduces membership to the pointwise commutation
equation.

The range is a block-diagonal commutant and need not be commutative. For this
reason, we do not call it a ``classical algebra''; the classical layer is the
span of Section~\ref{sec:average}. The library has no theorem identifying
the span with the center of the commutant, and no rank-one
classical-representation theorem.

\begin{theorem}[Bundled retraction and exact range]
\label{thm:pinching-range}
For every projective partition $\pi$, the map $\pinch_\pi$ is complex linear,
unital, positive, trace preserving, and idempotent. It lands in $N_\pi$ and
fixes $N_\pi$ pointwise. Consequently
\[
 \pinch_\pi(X)=X \Longleftrightarrow X\in N_\pi,
 \qquad
 \operatorname{range}(\pinch_\pi)=N_\pi
\]
where the second equality is the equality of the linear-map range and the
underlying submodule of the bundled commutant.
\end{theorem}

The corresponding declarations are
\texttt{partitionPinching\_unital},
\texttt{partitionPinching\_posSemidef},
\texttt{trace\_partitionPinching},
\texttt{partitionPinching\_idem},
\texttt{partitionPinching\_eq\_self\_iff\_mem\_commutant}, and
\texttt{range\_partitionPinchingLinearMap}. State preservation is exposed by
\texttt{partitionPinching\_isState}.

\begin{theorem}[Bimodule and Hilbert--Schmidt structure]
\label{thm:pinching-geometry}
If $A,B\in N_\pi$, then
\[
 \pinch_\pi(AXB)=A\pinch_\pi(X)B.
\]
Pinching is also selfadjoint for the trace pairing and satisfies the
Pythagorean identity
\[
 \Tr(X^*X)=
 \Tr(\pinch_\pi(X)^*\pinch_\pi(X))+
 \Tr((X-\pinch_\pi(X))^*(X-\pinch_\pi(X))),
\]
and hence contracts the squared Hilbert--Schmidt quantity. A
commutant-valued complex-linear map with the same trace pairings against every
commutant element equals
\texttt{partitionPinching\allowbreak LinearMap}.
\end{theorem}

The corresponding declarations are
\texttt{partitionPinching\_bimodule},
\texttt{trace\_conjTranspose\_partitionPinching\_mul},
\texttt{trace\_partitionPinching\_pythagoras},
\texttt{trace\_partitionPinching\_contraction},
\texttt{trace\_partitionPinching\_mul\_commutant}, and
\texttt{partitionPinchingLinearMap\_unique}. Together they make the map a
positive, unital, trace-preserving linear projection with a bimodule law and
a trace-duality characterization; these are the identities that
characterize the trace-preserving conditional expectation onto a subalgebra
in the operator-algebra literature \cite{umegaki1954,takesaki1972}. The
library proves the identities directly for this map. This core module does
not instantiate an operator-algebra conditional-expectation interface. The
companion \texttt{Dynamics/ChoiCPTP.lean} separately proves its custom
finite-matrix \texttt{IsCPTP} predicate from the same Kraus and trace
identities.

\subsection{L\"uders naturality}

If $P\in N_\pi$, compression commutes with pinching. After normalization,
trace preservation against commutant elements gives
\[
 \pinch_\pi(\lued_P(\rho))=
 \lued_P(\pinch_\pi(\rho)).
\]
The raw formula is
\texttt{partitionPinching\_luedersUpdate}. The theorem
\texttt{partitionPinching\_luedersStateUpdate} states the same naturality for
\texttt{StateMatrix n} and \texttt{ProjectionEvent n}: the input and output
are states, and the nonzero-weight proof is transported by
\texttt{bornWeight\_partitionPinching}. This checks the same operation through
the bundled state and event interfaces.

\section{Random-unitary pinching and the entropy support boundary}
\label{sec:spectral-boundary}

The first B9 majorization precursor has an exact finite representation. For
each sign assignment \(s:\{0,\ldots,k-1\}\to\{\pm1\}\), set
\[
 U_s=\sum_i s(i)P_i .
\]
Orthogonality and completeness make every \(U_s\) a selfadjoint unitary. The
independent sign characters obey
\(\sum_s s(i)s(j)=2^k\delta_{ij}\), hence
\[
 \frac1{2^k}\sum_s U_s X U_s^*=\sum_iP_iXP_i=\pinch_\pi(X).
\]
The real weights are strictly positive and sum to one, so this is a literal
random-unitary representation. A two-coordinate control proves that replacing
the independent sign family by only the global pair \(\{I,-I\}\) leaves an
off-diagonal matrix unchanged and is not pinching. These statements are
checked in \texttt{RecordMajorization.lean}. They are an algebraic precursor;
they do not prove spectral majorization or entropy monotonicity.

The support layer cannot be bypassed by totalizing the logarithm. In Lean,
\(\log 0=0\); continuous functional calculus therefore sends the logarithm
of every projection to the zero matrix, since a projection's spectrum lies
in \(\{0,1\}\). Consequently the raw finite expression
\[
 \operatorname{Re}\Tr\!\left[\rho(\log\rho-\log\sigma)\right]
\]
is zero for the two orthogonal rank-one coordinate projections. Both are
density matrices, but the first support is not contained in the second.
Thus any extended-real divergence that returns \(+\infty\) on support failure
cannot equal this totalized raw formula. The density, support, zero-value, and
inequivalence receipts are checked in
\texttt{SpectralEntropyBoundary.lean}.

This is a fail-closed architectural result, not a no-go for Umegaki entropy.
A complete continuation must define the support-aware extended divergence and
then prove the pinching relative-entropy Pythagorean identity, spectral
majorization, constrained maximum-entropy formula, and the two-level
publicization information chain.

\section{Averaging on the partition span}
\label{sec:average}

The same partition carries a second canonical expectation. The
\emph{partition span}
\[
 D_\pi=\operatorname{span}\{P_i\}
\]
is bundled as \texttt{ProjectivePartition.span}. Its closure package is
proved by span induction.

\begin{proposition}[The commutative layer]
\label{prop:span}
$D_\pi$ contains every projector and the identity, and is closed under
multiplication and conjugate transposition
(\texttt{proj\_mem\_span}, \texttt{one\_mem\_span},
\texttt{mul\_mem\_span}, \texttt{conjTranspose\_mem\_span}). It is
commutative (\texttt{span\_mul\_comm}). It is contained in the commutant
(\texttt{span\_le\_commutant}), and every element of $D_\pi$ commutes with
every element of $N_\pi$
(\texttt{mul\_comm\_of\_mem\_span\_of\_mem\_commutant}): the span lies in the
center of the commutant.
\end{proposition}

Only the inclusion into the center is used here. The reverse inclusion is not
part of this module's API; zero projectors are removed in the separate active-label
coordinate equivalence. The averaging map is
\[
 \avg_\pi(X)=\sum_i\frac{\Tr(XP_i)}{\Tr(P_i)}\,P_i,
\]
bundled as \texttt{partitionAverageLinearMap}. Lean's zero-totalized inverse
makes the terms of zero projectors vanish, so no nonzero-block hypothesis
appears anywhere in the module: a zero projector contributes zero to both
sides of every identity below.

\begin{theorem}[Averaging expectation and exact range]
\label{thm:average}
For every projective partition $\pi$, the map $\avg_\pi$ is complex linear,
unital, positive, trace preserving, and idempotent. It lands in $D_\pi$,
fixes $D_\pi$ pointwise, and
\[
 \avg_\pi(X)=X \Longleftrightarrow X\in D_\pi,
 \qquad
 \operatorname{range}(\avg_\pi)=D_\pi.
\]
Against every $C\in D_\pi$ the average is invisible to the trace pairing,
$\Tr(\avg_\pi(X)\,C)=\Tr(XC)$, and any $D_\pi$-valued complex-linear map with
this property equals the bundled average.
\end{theorem}

The corresponding declarations are
\texttt{partitionAverage\_unital},
\texttt{partitionAverage\_posSemidef},
\texttt{trace\_partitionAverage},
\texttt{partitionAverage\_idem},
\texttt{partitionAverage\_eq\_self\_iff\_mem\_span},
\texttt{range\_partitionAverageLinearMap},
\texttt{trace\_partitionAverage\_mul\_of\_mem}, and
\texttt{partitionAverageLinearMap\_unique}; state preservation is
\texttt{partitionAverage\_isState}. The uniqueness proof mirrors the pinching
case: the difference lies in the star-closed span, is trace-orthogonal to its
own conjugate transpose, and vanishes by faithfulness of the trace.

\begin{theorem}[Two-level structure]
\label{thm:two-level}
The two expectations compose as a tower,
\[
 \avg_\pi\circ\pinch_\pi=\avg_\pi=\pinch_\pi\circ\avg_\pi ,
\]
(\texttt{partitionAverage\_partitionPinching},
\texttt{partitionPinching\_partitionAverage}). The average preserves the
Born statistics of the partition,
$\mu_{\avg_\pi(\rho)}(P_j)=\mu_\rho(P_j)$ for every $j$
(\texttt{bornWeight\_partitionAverage}). For a partition member $P_j$ of
nonzero weight, averaging commutes with L\"uders conditioning, and both
composites collapse to the normalized projector:
\[
 \avg_\pi(\lued_{P_j}(\rho))
 =\lued_{P_j}(\avg_\pi(\rho))
 =\Tr(P_j)^{-1}\,P_j
\]
(\texttt{partitionAverage\_luedersUpdate},
\texttt{partitionAverage\_luedersUpdate\_proj},
\texttt{luedersUpdate\_partitionAverage\_proj}).
\end{theorem}

The collapse identity is the finite-matrix form of classical conditioning:
on the commutative layer, conditioning on an observed partition member
replaces the averaged state by the indicator of that member, normalized. The
identity is stated for raw matrices; the only hypothesis is the nonzero
outcome weight.

The same partition also gives an exact operational quotient. Two matrices
$X$ and $Y$ are declared equivalent when

\[
 \Tr(XC)=\Tr(YC)
 \qquad\text{for every }C\text{ commuting with all }P_i.
\]

\begin{theorem}[Complete partition quotient]
\label{thm:partition-operational-quotient}
For every finite projective partition,
\[
 \bigl(\forall C\in N_\pi:\Tr(XC)=\Tr(YC)\bigr)
 \quad\Longleftrightarrow\quad
 \pinch_\pi(X)=\pinch_\pi(Y).
\]
Consequently, every matrix $D$ with $\pinch_\pi(D)=0$ is invisible to all
trace tests from the sector-preserving commutant. In particular, each cross-sector corner
$P_iDP_j$ with $i\ne j$ lies in that kernel.
\end{theorem}

\begin{proof}
Trace duality gives the forward statistics of the pinched matrices. For the
converse, apply the uniqueness theorem for the commutant-valued trace-dual
projection. The cross-sector statement follows from orthogonality of the
partition projectors. The complete equivalence, kernel, and corner statements
are machine-checked in \texttt{Superselection.lean}.
\end{proof}

The theorem is relative to a supplied partition. It neither constructs a
physical sector decomposition nor identifies the commutant with a laboratory
readout algebra.

For every supplied projective partition, averaging maps surjectively onto the
commutative partition-span algebra. It factors through block pinching, and
every value commutes with the partition commutant. A one-block rank-two
control keeps the maps distinct: pinching retains a within-block coherence
which averaging removes. The finite matter interface bundles the same map as
a typed public-record adaptor. These statements are machine-checked in
\texttt{B10EdgeCenterAction.lean}. The interface adds no edge object,
edge-to-partition identification, source selection, or physical detector
interpretation to the event-algebra theorems.

\section{The bundled state expectation}
\label{sec:functional}

For a matrix $\rho$, \texttt{stateExpectationLinearMap rho} bundles
$M\mapsto\Tr(\rho M)$ as a complex-linear functional. If $\rho$ is a state it
maps identity to one. If $\rho$ and $M$ are positive semidefinite, then the
expectation is nonnegative. The proof of
\texttt{expectation\_nonneg} uses a finite spectral decomposition of $M$,
cycles the trace, and reduces the result to a sum of products of nonnegative
diagonal entries; Section~\ref{sec:mathlib} explains why this elementary
route was chosen. Additivity and homogeneity come from the bundled
\texttt{LinearMap} interface and are not repeated as pointwise lemmas.

\section{Finite-state Robertson uncertainty}
\label{sec:robertson}

Let $\rho$ be a supplied density matrix and let $A$ and $B$ be supplied
Hermitian matrices. Write
\[
 A_0=A-\Tr(\rho A)\one,
 \qquad
 B_0=B-\Tr(\rho B)\one,
\]
and define
\[
 (\Delta_\rho A)^2=\Re\Tr(\rho A_0^2),
 \qquad
 c_\rho(A,B)=2\Im\Tr(\rho A_0B_0).
\]
Hermiticity makes both subtracted expectations real. The positive matrix
$\rho$ induces the seminormed pairing
\[
 \langle X,Y\rangle_\rho=\Tr(Y\rho X^*).
\]
It may be degenerate when $\rho$ is not faithful, which is why the formal
proof uses Mathlib's positive-semidefinite matrix seminorm rather than
silently assuming a positive-definite state.

\begin{theorem}[Supplied-state Robertson inequality]
\label{thm:finite-robertson}
For every supplied finite state and pair of Hermitian observables,
\[
 \frac{c_\rho(A,B)^2}{4}
 \le (\Delta_\rho A)^2(\Delta_\rho B)^2.
\]
The same finite pairing also gives
\[
 \bigl(c_\rho(A,B):\mathbb C\bigr)=
 -i\Tr\bigl(\rho(AB-BA)\bigr).
\]
In particular, the ordinary-commutator form is
\[
 \frac{\left|\Tr\bigl(\rho(AB-BA)\bigr)\right|^2}{4}
 \le (\Delta_\rho A)^2(\Delta_\rho B)^2.
\]
A zero variance on either side therefore forces the commutator readout to
vanish.
\end{theorem}

\begin{proof}
Cauchy--Schwarz bounds the squared modulus of
$\langle A_0,B_0\rangle_\rho$. Its imaginary part is no larger. Trace
cyclicity identifies the pairing with $\Tr(\rho A_0B_0)$, and scalar centring
does not change the commutator. The complex identity, the centered reduction,
and the ordinary-commutator inequality are machine-checked in
\texttt{Robertson.lean}.
\end{proof}

The exact two-level control uses $\rho=|{+z}\rangle\langle{+z}|$.
Pauli $X$ and $Y$ do not commute, have unit variance, and saturate the bound
with $c_\rho(X,Y)=2$. Pauli $Z$ and $X$ also do not commute, while $Z$ has
zero variance and the state commutator readout vanishes. The theorem is an
uncertainty statement for supplied finite inputs. It selects no state,
observable, Hamiltonian, or physical instrument.

\section{CHSH interoperability as a client}
\label{sec:chsh}

For selfadjoint involutions $a_0,a_1,b_0,b_1$ with cross commutation, set
\[
 S=a_0b_0+a_0b_1+a_1b_0-a_1b_1.
\]
The development first proves, in a bare ring,
\[
 S^2=4\one-(a_0a_1-a_1a_0)(b_0b_1-b_1b_0)
\]
(\texttt{chsh\_mul\_self}). In a nontrivial unital C*-ring this yields the
operator-norm bound $\norm{S}\le2\sqrt2$
(\texttt{tsirelson\_bound}), the norm form of Tsirelson's inequality
\cite{tsirelson1980} for the CHSH combination \cite{chsh1969}. The analytic
half is short: selfadjoint involutions have norm one by the C*-identity
(\texttt{norm\_eq\_one\_of\_selfAdjoint\_involution}), commutators of
norm-one elements have norm at most two
(\texttt{norm\_commutator\_le\_two}), and the square identity converts these
bounds into $\norm{S}^2\le8$. The precise Mathlib structure is a
\texttt{NormedRing}, \texttt{StarRing}, \texttt{CStarRing}, and
\texttt{Nontrivial}. The result is stated for a unital C*-ring.

The theorem \texttt{tsirelson\_bound\_of\_isCHSHTuple} accepts Mathlib's
\texttt{IsCHSHTuple}. The matrix instantiation,
\texttt{matrix\_tsirelson\_bound}, activates Mathlib's scoped operator
norm. The declaration \texttt{matrix\_tsirelson\_bound\_of\_events} accepts four
projection events plus the four cross-commutation equations, constructs the
dichotomic observables $\one-2P$, and returns the norm bound. This gives the
event layer a compiled path into the CHSH theorem.

Mathlib's \texttt{tsirelson\_inequality} is an order inequality in a
star-ordered real algebra \cite{mathlibCHSH2026}. The theorems in this
artifact are an operator-norm bound, its state-level corollary below, and
the exact equality witness of Section~\ref{subsec:saturation}.

\subsection{The state-level corollary}
\label{subsec:statechsh}

The bridge from norms to states is the checked bound
\[
 \lvert\Tr(\rho M)\rvert\le\norm{M}
\]
for every state $\rho$ and observable $M$
(\texttt{norm\_expectation\_le\_l2\_opNorm}). The proof diagonalizes the
state, cycles the trace to $\Tr\!\bigl(D\,(V^*MV)\bigr)$, dominates every
diagonal entry of the conjugated observable by the operator norm
(\texttt{norm\_diag\_entry\_le\_l2\_opNorm}, by testing against a standard
basis vector), uses that unitary conjugation does not increase the norm
(\texttt{norm\_eq\_one\_of\_mem\_unitaryGroup} with submultiplicativity), and
sums the nonnegative eigenvalues to one. Composing with the event-level norm
bound gives the state-level CHSH corollary
(\texttt{matrix\_state\_tsirelson\_bound\_of\_events}): for a state and four
projection events with cross-party commutation,
\[
 \bigl\lvert\Tr(\rho S)\bigr\rvert\le2\sqrt2
\]
for the CHSH combination $S$ of the dichotomic observables $\one-2P$. The
pairing $\Tr(\rho S)$ is real by \texttt{star\_bornWeight}, since $S$ is
selfadjoint, so the modulus bound is a two-sided real bound.
Section~\ref{subsec:saturation} exhibits a declared witness attaining the
bound.

\subsection{Exact saturation on the slot-locality interface}
\label{subsec:saturation}

The module \texttt{TsirelsonSaturation.lean} supplies the attainment
witness. Its carrier objects come from the wider library. A supplied
bipartite slot split of the record algebra, committed as the interface
\texttt{SlotLocalityInterface}, which carries a norm-form Tsirelson bound
through the theorem \texttt{slot\_locality\_receipts}. The module inhabits
that interface at two-dimensional slots with the singleton identity Kraus
family (\texttt{bellSlotInterface}). The declared state is the normalized
maximally entangled state on the four-dimensional product carrier,
transported to $\mathrm{Fin}\,4$ along a fixed index equivalence so that
the typed state and expectation vocabulary of
Sections~\ref{sec:setting} and \ref{sec:functional} applies
(\texttt{bellState\_isState}). The four declared events are the spectral
projections of the two Pauli directions on the left slot and of the two
quarter-rotated directions on the right slot; each carries a checked event
certificate, and the four cross-party commutations are checked on the
product carrier.

\begin{theorem}[Exact Tsirelson saturation]
\label{thm:saturation}
The CHSH expectation of the declared state $\rho_{\mathrm{Bell}}$ against
the combination $S$ of the dichotomic observables $\one-2P$ of the four
declared events equals the upper bound exactly:
\[
 \Tr(\rho_{\mathrm{Bell}}\,S)=2\sqrt2 .
\]
\end{theorem}

The equality is \texttt{chsh\_saturation}; its modulus form is
\texttt{chsh\_saturation\_norm}. The packaged statement,
\texttt{tsirelson\_saturation\_receipt}, is one conjunction over one
witness: the committed norm bound of the interface cited at the witness,
the state-level bound of Section~\ref{subsec:statechsh}
(\texttt{matrix\_state\_tsirelson\_bound\_of\_events}) cited at the four
events for every state on the carrier, the state certificate with the
exact attainment in value and in modulus, and the identification of the
attaining observable with the slot-lifted CHSH operator of the interface
witness (\texttt{toFour\_chshProd\_eq}). The two upper bounds are cited
from the committed theorems, not re-proved.

The boundary is declared data. The saturating state and the four events
are declarations of the module; no theorem produces them from committed
source runs, and no physical regions, spacelike separation, or observer
instruments are constructed. Identification of the finite causal structure with an
observer-produced physical spacetime stays open, and the receipt consumes the
bipartite slot split as a supplied datum.

\subsection{The record-diagonal classical delimitation}
\label{subsec:delimitation}

The same module bounds the classical side. For convex weights $w$ and four
diagonal observables $a_0,a_1,b_0,b_1$ valued in $[-1,1]$ on any finite
carrier, the CHSH expectation obeys the classical bound
\[
 \Bigl|\sum_i w_i\,
 (a_{0i}b_{0i}+a_{0i}b_{1i}+a_{1i}b_{0i}-a_{1i}b_{1i})\Bigr|\le2
\]
(\texttt{diagonal\_chsh\_le\_two}); the matrix form, for a diagonal state
read against four diagonal readouts in the CHSH combination, is
\texttt{diagonal\_state\_chsh\_le\_two}. In particular, the counted
correlation state of the wider library, a record-diagonal state assembled
from the counted joint run of the committed observer pair $(86,247)$,
obeys the bound $2$ against every quadruple of diagonal $[-1,1]$-valued
readouts (\texttt{counted\_state\_diagonal\_chsh\_le\_two}). Record-basis
diagonality is therefore insufficient for the value $2\sqrt2$. This is a
delimitation of one fixed diagonal state with four readouts diagonal in the
same basis, not of arbitrary contextual classical experiments. The Bell
state and all four attaining settings are real matrices, so the proved gap is
entanglement, non-diagonal coherence, and noncommutativity; it does not isolate
the genuinely complex Pauli-$Y$ direction of the declared phase-sensitive
effect, and no theorem asserts that the committed phase lift is necessary or
sufficient for attainment.

The delimitation extends from one fixed counted state to the whole committed
production surface
(\texttt{SourceReachabilityDelimitation.lean}). The module defines the class
of states reachable from the three source-counted occupation laws under
every operation the development commits: the walk-step transports, the
internal slot exchange, the ambient anchorings, the marginalizations, the
regional and scalar conditional expectations, marginal products, and convex
mixtures. Every state in that class remains diagonal in the record basis
with nonnegative weights summing to one, and its CHSH value against any four
record-diagonal $[-1,1]$-valued readouts is at most $2$, strictly below the
attained $2\sqrt2$ of the declared witness. The witness state itself carries an
off-diagonal entry of one half, so no map that preserves record diagonality
reaches that particular witness from the class. This target-specific
obstruction does not show that every realization of the value $2\sqrt2$
requires non-diagonal preparation: readouts that are not slot-local with
respect to the declared split are outside the bound, and the slot-local
separability theorem for arbitrary settings is supplied by the companion
module described next. The result is a statement about the committed
operation class and the declared Bell-state witness; every conceivable
architecture extension is outside its scope. Like \texttt{Dynamics/ChoiCPTP.lean},
the module is a companion in the wider library, outside this paper's
fifteen-module count and scale table.

The companion module \texttt{ProductSplitSeparability.lean} lifts the jointly-diagonal
restriction on the declared slot split. A product-basis diagonal state with a probability
diagonal is the explicit convex combination of the point states $e_{aa}\otimes e_{bb}$
weighted by its diagonal entries (\texttt{productDiagonal\_separable}), and against any
four Hermitian settings in the Loewner unit interval, two on each slot and with no
commutation assumed inside a slot, its slot-local CHSH trace is real with modulus at most
$2$ (\texttt{productDiagonal\_slotLocal\_chsh\_le\_two}). Every reachable state on the
two committed pair carriers is therefore separable across the split and obeys the bound
against all slot-local unit-interval readouts. Slot membership is the hypothesis that carries
the bound: the Bell
settings conjugated by $\mathrm{CNOT}\,(H\otimes1)$ commute pairwise across the wings, lie
in the unit interval, and reach exactly $2\sqrt2$ on the product-diagonal state
$|00\rangle\langle00|$, with one of them in neither slot
(\texttt{crossWing\_commutation\_not\_sufficient}). A slot-local value above $2$ on a state certifies
off-diagonal record coherence; regions and source preparation stay with the open premises.
Like the reachability module, it is a companion in the wider library, outside the
fifteen-module count and scale table.

\section{An ideal static phase-POVM count fit}
\label{sec:phase-instrument}

The module \texttt{OperationalPhaseInstrument.lean} makes the interface
boundary machine-visible. Its structure
\texttt{IdealPhasePOVMCountModel} is an ideal static model: one matrix state,
the committed web effects plus one designated phase effect, exact binary POVM
certificates, integer count literals with the diagonal pair fixed to
$(111,68)$, and exact frequency-equals-Born equations. It contains no source
operation, completely-positive outcome maps, trace-preserving summed channel,
postmeasurement state, readback implementation, producer identity, receipt
binding, or statistical validation rule.

Every ideal model yields fixed-trace identification
(\texttt{ideal\_phase\_model\_completes\_tomography}), and its state is the
unique state fitting all exact equations
(\texttt{preparation\_identified\_by\_exact\_fit\_data}). The construction
\texttt{modelOfIdealPhaseFitData} is equivalent only to the existence of
\texttt{IdealPhaseFitData} (\texttt{hasIdealPhaseFit\_iff\_data}); this is an
algebraic fit statement rather than a theorem establishing the
phase-sensitive effect.

Two incompleteness controls are exact. Two certified distinct states carry identical Born
data on every diagonal effect (\texttt{diagonal\_context\_insufficient}),
so the diagonal context alone identifies no state; the committed rotated
context separates that pair; and no diagonal effect and no complexified
real effect satisfies the phase-sensitivity clause
(\texttt{phase\_effect\_not\_diagonal},
\texttt{phase\_effect\_genuinely\_complex}). Hence a complete static family
needs non-diagonal content and some separator outside the complexified-real
closure (\texttt{phase\_context\_necessary}). The theorem does not force the
particular rotated effect, the designated phase effect, their order, or an
operational implementation.

The non-discharge boundary is exact. The declaration
\texttt{syntheticIdealPhaseFitData} chooses the rational diagonal state with
weights $111/179$ and $68/179$, rotated counts $(315,401)$, and phase counts
$(1,1)$. The theorem \texttt{synthetic\_ideal\_phase\_fit} proves that these
chosen literals inhabit \texttt{HasIdealPhaseFit} without any producer, run,
operation, or receipt input. Exact finite frequency equality is also not a
general validation criterion for sampled data. Therefore ideal fit cannot
certify custody or establish the phase-sensitive effect. A genuine target needs
source-attached CP outcome maps, a trace-preserving summed channel,
operation/readback and common-preparation semantics, producer-bound receipts,
and a preregistered statistical validation rule (or an explicitly exhaustive
deterministic semantics). Operational additivity and composition remain a
separate cross-context additivity obligation.

\subsection{Exact determination boundary and the integer receipt window}
\label{sec:phase-determination}

The module \texttt{PhaseInstrumentDetermination.lean} sharpens the boundary to
an exact reduction. The phase-sensitivity clause has a coordinate
characterization: for a Hermitian effect $E$ the Born-weight gap between the
two Pauli-Y states equals $-2\,\mathrm{Im}\,E_{01}$
(\texttt{bornWeight\_rhoY\_sub\_of\_isHermitian}), so the clause holds exactly
when $\mathrm{Im}\,E_{01} \neq 0$
(\texttt{phase\_sensitivity\_iff\_offdiag\_im}). The diagonal and
complexified-real exclusions are the $\mathrm{Im}\,E_{01} = 0$ instances, and
conjugating a committed web projector by any real matrix, gauge images and
permutations included, is excluded in one statement
(\texttt{phase\_effect\_ne\_real\_conjugated\_web}).

The committed run literals pin the model state. Every inhabitant satisfies
$\rho_{00} = 111/179$, $\rho_{11} = 68/179$, and
$\rho_{10} = \overline{\rho_{01}}$ (\texttt{prep\_coordinates}), so one
complex off-diagonal coordinate carries the entire remaining state content.
Every count frequency of every inhabitant, in every instrument context, is an
exact affine function of that coordinate, with real coefficients determined
by the context's effect entries (\texttt{frequency\_affine}). Over the
committed core the phase frequency equals
$1/2 - \mathrm{Im}\,\rho_{01}$ (\texttt{phase\_frequency\_eq}) and the rotated
frequency equals $315/716 - (\sqrt{3}/2)\,\mathrm{Re}\,\rho_{01}$
(\texttt{rotated\_frequency\_eq}); two committed-core inhabitants have equal
states exactly when all count frequencies agree
(\texttt{prep\_eq\_iff\_frequencies\_eq}). Positivity bounds the coordinate by
$|\rho_{01}|^2 \le 7548/32041$ (\texttt{prep\_offdiag\_normSq\_le}), which
forces the decidable integer window
$32041\,(a-b)^2 \le 30192\,(a+b)^2$ on the phase counts $(a,b)$ of every
committed-core inhabitant (\texttt{phase\_counts\_receipt\_window}). Both
phase outcomes of every committed-core inhabitant carry positive mass, and
the minimal phase count mass over committed-core inhabitants is exactly two,
attained by the synthetic inhabitant
(\texttt{minimal\_phase\_count\_mass}).

The reduction states the phase residue exactly: for a committed-core
inhabitant the state diagonal, the affine dependence of every count frequency
on the one free coordinate, and the admissible count window are forced by the
committed diagonal literals and the committed effects. Two operational requirements are not encoded by this static type: a
representation of the measurement by completely positive outcome maps with a
trace-preserving summed channel and compatible effects, and a source-produced
common preparation that selects the context and produces the public outcomes. The window is a
necessary arithmetic condition only and carries no provenance information,
since the synthetic inhabitant meets it without a source run.

\subsection{Static conformance fixture for the phase-effect layer}
\label{sec:phase-attainment}

The committed exact lift
$I/2 - (2\sqrt{3}/3)\,i\,(QP - PQ)$, equal to the Pauli $+Y$ projector, is taken
as a declared phase-sensitive effect. It is not a state transition or an
instrument. The module \path{OperationalPhaseAttainment.lean} consumes that
declaration and inhabits \texttt{IdealPhasePOVMCountModel} over the committed
core with static integer literals. A deterministic calculator expands the
Born table of the declared effects on the declared matrix
$\mathrm{diag}(111/179, 68/179)$ in exact arithmetic over
$\mathbb{Q}(\sqrt{3}, i)$: each context pair is the exact Born weight of the
context effect scaled to the least positive integer multiple of the reference
mass $179$ clearing its denominator, giving
$(111, 68)$ at mass $179$ in the diagonal and record-conjugate contexts,
$(315, 401)$ at mass $716$ in the rotated contexts, and $(179, 179)$ at mass
$358$ in the phase context. These are generated expected-frequency
numerators, not observations or validation. No sampling or floating-point
step occurs. An independent verifier sharing no arithmetic code rebuilds the table
from the pinned payload and demands exact agreement with every receipt
field, run-mass-multiple minimality and the integer window included. The
composed receipt (\texttt{operationalPhaseAttainment\_receipt}) derives,
from one antecedent bundle pinned to the declared effect and matrix, the
conjunction: committed core, declared effect in the phase slot, vanishing
off-diagonal matrix entry, committed diagonal literals, the stored Born-fit
equations and their cross-multiplied scaling consequences, positive integer
masses, and phase counts inside the registered window. The Lean theorem does
not prove leastness, source reachability of this two-dimensional matrix, or
receipt provenance; the external verifier checks the deterministic
least-denominator arithmetic separately. The instrument representation, the
source-produced preparation, and the independent cross-context additivity
premise all stay open.

The module \texttt{LuedersPhaseInstrument} types the channel clauses as a
structure \texttt{PhaseInstrument} on the eight committed contexts: outcome maps
that are completely positive, trace-nonincreasing on positive-semidefinite
inputs, summing to a trace-preserving channel per context, and inducing the
committed effect table. The L\"uders maps $X \mapsto EXE$ inhabit it
(\texttt{luedersPhaseInstrument}) with a singleton Kraus family, and the
normalized post-measurement state is certain of its effect
(\texttt{luedersPhaseInstrument\_repeatable}). On $\mathrm{diag}(111/179, 68/179)$
the outcome traces reproduce the fixture frequencies, with $1/2$ in the phase
context, whose post-measurement state is the declared $+Y$ projector itself;
the phase and diagonal values are also obtained by direct matrix computation
through the outcome maps, with no use of the stored Born-fit equations. A
swap-twisted instrument meets every channel clause with the same induced effects,
differs on the record entry, and fails repeatability
(\texttt{effect\_table\_does\_not\_determine\_instrument}). The L\"uders choice is a
declared selection, no source implements it, and the instrument representation,
the source-produced preparation, and cross-context additivity stay open. The module is a companion in the wider library, outside the
fifteen-module count and scale table.

The companion module \texttt{SourcePhaseSelection} asks whether the phase
slot of that instrument can be produced from source fields. Its generator
reads two inputs only, the six exact rows of the two-dimensional
irreducible representation and the diagonal record projector, forms the
orbit of the projector, enumerates the twelve noncommuting orbit pairs, and
normalizes each commutator by its own off-diagonal magnitude into a rank-one
projector. The generated effect values are exactly the two Pauli-$Y$
projectors (\texttt{generated\_effect\_values\_exact}); the first pair in
lexicographic index order selects the $+Y$ projector, which equals the
declared lift and the instrument's phase slot
(\texttt{source\_selected\_generated\_effect\_eq\_sourcePhaseLift},
\texttt{source\_selected\_generated\_effect\_eq\_declared\_lueders\_effect}).
Eight pairs give $+Y$ and four give $-Y$, so the theorem is effect-value
equality rather than witness uniqueness, and the orientation is fixed by two
producer conventions, the commutator order and the pair index: reversing
either selects the transpose effect
(\texttt{effectMatrix\_negative\_eq\_transpose\_positive}). What the source
fixes without convention is the unordered pair of $Y$ projectors. The
module's enabled-domain relation reproduces the canonical $36/12$ split
between state-changing L\"uders updates and equality stutters
(\texttt{changingEnabledCellCensus},
\texttt{alreadyOperationEnabledCellCensus}). It constructs no instrument
and supplies no preparation, outcome, readback, or provenance; the
phase-sensitive-effect declaration keeps the orientation choice. Like the
reachability module, it is a companion in the wider library, outside the
fifteen-module count and scale table.

The companion module
\texttt{SourcePhaseInstrumentOutcomeBridge} joins those generated source
semantics to the already declared L\"uders phase instrument. It maps every
generated $+Y$ event to the declared phase-outcome index $0$ and every
generated $-Y$ event to index $1$, then proves exact equality of the generated
effect with that entry of the committed effect pair
(\texttt{generatedEffect\_eq\_committedPhaseOutcome}). All five generated
state matrices are positive semidefinite with trace one
(\texttt{stateMatrix\_isState}). For each of the 48
enabled source steps the corresponding Born weight is nonzero, and the
normalized output of the indexed declared L\"uders map equals the generated
semantic result matrix
(\texttt{sourceStep\_normalized\_lueders\_outcome\_eq\_result}). This is a
matrix-level compatibility theorem. The index is not a recorded public
outcome, and the result neither source-selects nor implements the L\"uders
instrument, nor supplies one common cross-context preparation, readback, run
binding, provenance, or custody. The clauses are conjoined in
\texttt{sourcePhaseInstrumentOutcomeBridge\_receipt}.

The companion module \texttt{SourcePhaseCommonPreparationHull} identifies
the exact common-support surface of the generated phase relation. Exactly the
three real orbit states admit all twelve generated phase events
(\texttt{hasAllPhaseEvents\_iff}). Their normalized real mixture with weights
\[
  \frac{265}{537},\qquad \frac{136}{537},\qquad \frac{136}{537}
\]
is exactly the declared run state; this is theorem
\texttt{commonSupportMixture\_eq\_committedRunState}. Equality to that matrix
uniquely determines all three coefficients, without assuming normalization
(\texttt{commonSupportMixture\_coefficients\_unique}). Each common-support
state assigns weight $1/2$ to either generated Pauli-$Y$ orientation, and
every common-support state/event pair inherits the normalized declared
L\"uders-outcome equality. The declared run matrix is not any generated state,
no generated source step starts from it, and the result of every generated
phase event lacks all-event support. This is retrospective convex-hull
compatibility, not a source-produced or enabled common preparation. The
pointwise one-step cells supply no convex-lift, mixing, reset,
re-preparation, or sequential schedule. The twelve generated events are not
the eight public instrument contexts, and the theorem supplies no instrument
implementation, public outcome, readback, run binding, provenance, custody,
cross-context additivity, or physical attachment; PR-65 remains open. The
clauses are conjoined in
\texttt{sourcePhaseCommonPreparationHull\_receipt}.

The module \texttt{SourcePhaseBornWeightBoundary} proves the exact
single-shot identity
\[
  \operatorname{Tr}(\rho P_{+Y})=\frac12-\operatorname{Im}(\rho_{01})
\]
for every normalized two-dimensional state. Thus a non-half selected-effect
weight requires nonzero imaginary off-diagonal coherence; real off-diagonal
coherence alone is insufficient
(\texttt{state\_sourceSelected\_weight\_eq\_half\_sub\_im}). It then places
the selected effect and the enumerated source-reachable class behind an
explicit prospective preparation-bridge interface. If such a bridge maps
reachable matrices to states and preserves record diagonality, the selected
$+Y$ effect has Born weight exactly $1/2$ on every bridged preparation
(\texttt{reachable\_recordPreserving\_phase\_weight\_half}). On a
record-diagonal state the generated $-Y$ orientation and the diagonal
comparator $(1/2)I$ have the same weight. Constant diagonal and off-diagonal
bridges show that the preservation hypothesis is logically inhabitable and
load-bearing. Both are mathematical controls that ignore their inputs. The
result derives no source bridge or operation, common preparation, public
outcome, sequential instrument behavior, operational readback, run,
provenance, or custody.

\section{Working over Mathlib: scopes, gaps, and affordances}
\label{sec:mathlib}

This section records what the construction consumed from Mathlib, where the
library resisted, and which workarounds closed the gaps. The artifact
repeats the catalog in the notes file \texttt{MATHLIB\_NOTES.md}.

\subsection{Scoped instances}

The three scopes of Section~\ref{sec:setting} share a failure mode. The
partial order on $\C$ and the Loewner order are scoped instances; every
$0\le z$ goal and every positive-semidefiniteness statement over $\C$
elaborates only when the scope is open, and a missing scope surfaces as an
instance-resolution error far from its cause. \texttt{ComplexOrder} also
carries the scoped ordered-ring structure on $\C$ that supplies
multiplicativity of nonnegativity inside the complex order. There is no
global norm on Mathlib matrices; the operator norm, the
\texttt{NormedRing} structure, and the \texttt{CStarRing} instance arrive
only under \texttt{Matrix.Norms.L2Operator}, so the entire Tsirelson section
lives behind that scope. The \texttt{Nontrivial} instance for
$n\times n$ complex matrices with $n\ne0$ did not synthesize and is
constructed by hand, entrywise. The matrix unitary group is a submonoid
distinct from Mathlib's \texttt{unitary} bundle, so its norm-one property is
also proved by hand from the C*-identity.

\subsection{Gaps and workarounds}

\emph{Bilinear trace positivity.} Mathlib proves $0\le\Tr(A)$ for a single
positive-semidefinite matrix and covers sandwich forms $BAB^*$. The
bilinear statement $0\le\Tr(AB)$ for two positive-semidefinite factors is
absent. For Born weights the gap costs nothing, because
$\Tr(\rho P)=\Tr(P\rho P)$ for idempotent $P$. For
\texttt{expectation\_nonneg} the natural route, the C*-order equivalence
between nonnegativity and Gram factorization, fails to elaborate: two
continuous-functional-calculus instances are not synthesized for the matrix
algebra under its product topology. The proof takes the elementary route
instead, a spectral decomposition of the observable followed by a trace
cycle and termwise nonnegativity, in about a dozen lines.

\emph{The ring identity.} The square identity behind the Tsirelson bound is
pure noncommutative ring algebra under side relations: four involutivity
equations and four cross commutations. Lean has no noncommutative analogue
of \texttt{linear\_combination}, and \texttt{noncomm\_ring} does not use
hypotheses. The checked proof builds a small confluent rewrite system. Each
commutation hypothesis is recast as a universally quantified,
association-compatible form $b(ax)=a(bx)$, and each involution as
$a(ax)=x$; applied through \texttt{simp only}, these rules move every $b$
past every $a$ in right-associated words and cancel adjacent equal letters.
Every rewrite strictly decreases the number of letter pairs out of normal
order, so the system terminates, and \texttt{abel} closes the remaining
additive goal.

\emph{Higher-order motives.} Two steps needed their motives spelled out.
The binary span-induction principle behind the closure package of
Proposition~\ref{prop:span} does not infer its motive from the goal; the
predicate is passed explicitly. Rewriting with the unit-trace equation
inside a hypothesis whose proof term mentions the bundled state predicate
fails the motive typecheck, because the state predicate itself contains the
trace being abstracted; a term-level chain through injectivity of the real
embedding replaces the rewrite.

\emph{Rewriting discipline.} Unfolding the definition of the Born weight
rewrites the first syntactic occurrence, which in normalized expressions is
often the normalizing scalar $\mu_\rho(P)^{-1}$ instead of the intended trace
argument. The affected proofs isolate each trace identity as a standalone
step with explicit arguments to the trace commutation and cycling lemmas.

\subsection{Affordances}

The \texttt{Matrix.PosSemidef} API is complete enough that the development
never unfolds the definition of positive semidefiniteness: sandwich closure,
sums, scalar multiples, diagonal nonnegativity, trace nonnegativity, and the
Gram construction cover every need, and the scalar lemma accepts a complex
scalar that is nonnegative in the complex order, which makes the L\"uders
normalization a single application. Trace faithfulness is available in two
forms: the conjugate-transpose criterion powers the uniqueness theorems of
both expectations, and the zero-trace criterion together with the
Cauchy--Schwarz property of positive matrices powers the support theorem of
Section~\ref{sec:lueders}. On the analytic side, the C*-identity, the
norm-one property of the unit, and submultiplicativity reduce the norm half
of the Tsirelson bound to about forty lines over the four-typeclass
signature, and the matrix case is a two-line instantiation behind the norm
scope. The state-expectation bound of Section~\ref{subsec:statechsh} rests
on the scoped norm being definitionally the Euclidean operator norm: the
type synonym is identity-transparent, so a \texttt{show} converts
Euclidean-space statements to raw matrix-vector algebra, and the
basis-vector, submultiplicativity, and eigenvalue-trace lemmas finish the
proof in about sixty lines. Mathlib's CHSH file supplied the
\texttt{IsCHSHTuple} bundle, so interoperability with the order-form theorem
cost one conversion theorem.

\section{Artifact evaluation and proof audit}
\label{sec:evaluation}

The source is a Lake project. Its toolchain pins
\texttt{leanprover/lean4:v4.29.1}; its manifest pins Mathlib commit
\texttt{5e932f97dd25535344f80f9dd8da3aab83df0fe6}. Build the library with
\begin{center}
\texttt{lake build EventAlgebra}.
\end{center}
The canonical Lean modules are publicly available in the Observer Patch
Holography repository \cite{observerPatchHolographyRepo}. The exact source
snapshot is the repository revision containing this manuscript and its built
PDF; the manuscript source accompanies the submission.

Table~\ref{tab:scale} describes the fifteen modules discussed here and the
umbrella root. Other modules imported by the broader root are outside this
paper's count. The comparison in Section~\ref{sec:comparison} concerns
interfaces, not counts.

\begin{table}[ht]
\centering
\small
\begin{tabular}{@{}lrr@{}}
\toprule
Module & Lines & Audited declarations \\
\midrule
\texttt{Basic} & 302 & 25 \\
\texttt{Lueders} & 294 & 14 \\
\texttt{PartitionPinching} & 474 & 27 \\
\texttt{PartitionAverage} & 513 & 30 \\
\texttt{StateExpectation} & 105 & 4 \\
\texttt{Robertson} & 416 & 20 \\
\texttt{Superselection} & 197 & 9 \\
\texttt{RecordMajorization} & 322 & 6 \\
\texttt{SpectralEntropyBoundary} & 161 & 5 \\
\texttt{Tsirelson} & 275 & 8 \\
\texttt{ExpectationBound} & 174 & 4 \\
\texttt{TsirelsonSaturation} & 834 & 39 \\
\texttt{OperationalPhaseInstrument} & 878 & 40 \\
\texttt{PhaseInstrumentDetermination} & 503 & 24 \\
\texttt{OperationalPhaseAttainment} & 383 & 22 \\
\texttt{EventAlgebra} (root) & 147 & 0 \\
\midrule
Total & 5978 & 277 \\
\bottomrule
\end{tabular}
\caption{Source lines and per-module counts of declarations covered by the
generated axiom audit.}
\label{tab:scale}
\end{table}

Every public result discussed in the paper appears under its Lean name. The
modules end with \texttt{\#print axioms} commands for the principal
declarations, and a clean build emits the generated output. The 277 audited
declarations report only subsets of propositional extensionality, classical
choice, and quotient soundness, the standard Lean/Mathlib base. No declaration
reports \texttt{sorryAx}, and a source scan finds no \texttt{sorry}
placeholder in these modules.

The submitted API uses every orthogonality hypothesis in its orthogonal-sum
theorem. The sequential-update theorem has no unused event premise. State
update is typed and excludes dimension zero. The fixed-point equivalence has
no nonzero-weight guard, and the averaging module carries no nonzero-block
hypotheses. The commutant, span, and both expectations are bundled, and the
CHSH adapters are clients of the event definitions.

The repository contains the toolchain, Lake manifest, source modules, build
command, in-source axiom queries, and a Creative Commons BY-NC-SA 4.0 license.
No persistent archive identifier is claimed.

\section{Feature-level comparison}
\label{sec:comparison}

Table~\ref{tab:comparison} compares compiled features. It does not treat a
count of elementary lemmas as a research contribution.

\begin{table}[ht]
\centering
\scriptsize
\begin{tabular}{@{}L{2.1cm}L{4.25cm}L{6.0cm}@{}}
\toprule
Feature & Closest checked infrastructure & This artifact \\
\midrule
Projective measurement & Isabelle/HOL projective-measurement developments;
state/measurement interfaces in current Lean libraries & finite matrix event
predicate, typed event/state subtypes, closure and trace-pairing lemmas; treated
as prerequisite; no priority claim \\
Pinching & Physlib spectral pinching is a state-selected bundled CPTP map with
fixed-point and geometric results & arbitrary supplied projective partition;
bundled linear map and commutant; companion custom-\texttt{IsCPTP} proof; no
Physlib correspondence theorem \\
Retraction structure & operator-algebra conditional expectations and channel
interfaces & exact range/fixed points, positive/unital/trace-preserving,
bimodule, Hilbert--Schmidt geometry, map-level trace-duality uniqueness \\
Random-unitary and entropy boundary & random-unitary dephasing and
support-aware Umegaki divergence in finite quantum information & exact
independent-sign average for arbitrary supplied partitions; a support-failure
countermodel showing that totalized matrix logarithms cannot encode the
extended divergence \\
Commutative layer & conditional expectations onto commutative subalgebras in
operator-algebra practice & bundled averaging expectation onto the partition
span, exact range and uniqueness, tower laws with pinching, Born-statistics
preservation, classical-conditioning collapse \\
L\"uders update & Isabelle projective measurement and general channel
infrastructure & raw formula plus typed partial state API, support theorem,
unguarded state fixed-point equivalence, typed naturality with
arbitrary-partition pinching \\
CHSH/\allowbreak Tsirelson & Mathlib order theorem; Isabelle upper bound; Lean CHSH
rigidity & complementary norm theorem, event-level client, state-expectation
corollary, and an exact saturation witness on declared data; fixed
record-diagonal states with jointly diagonal readouts capped at the classical
bound \\
Reproducibility & public CI and archived releases are common best practice &
pinned Lean/Mathlib, an exact repository snapshot, clean build instructions,
and declaration-level axiom queries; no DOI claim \\
\bottomrule
\end{tabular}
\caption{Direct comparison with the closest formal infrastructure.}
\label{tab:comparison}
\end{table}

Physlib provides more channel structure: its spectral pinching is a
\texttt{CPTPMap} \cite{physlibPinching2026}. This artifact accepts a supplied
partition that need not come from the spectrum of a state. The two APIs
therefore have different inputs, and no correspondence theorem connects them.

Lean-Quantum and Lean-QIT cover a larger part of quantum information
\cite{kasauraEtAl2026LeanQuantum,zhuEtAl2026LeanQIT}. Here, the scope is
limited to connections among projections, commutants, spans, the two
expectations, selective update, and Mathlib's CHSH interface.

\section{Limitations and conclusion}
\label{sec:conclusion}

The library connects projection events, Born weights, a typed L\"uders
update, arbitrary projective partitions, the two expectations of a partition,
and event-level CHSH bounds in norm and state form, with an exact saturation
witness and a record-diagonal classical delimitation on the same event API. The pinching has the
bundled commutant as its exact range and the average has the bundled span;
each map is the unique trace-dual projection onto its range, the two compose
as a tower, and both interact with L\"uders conditioning by checked
theorems. The supplied-state Robertson bound, the independent-sign
representation, the logarithm-support boundary, and the partition-relative
operational quotient use the same finite matrix API. Each formal statement
in the paper points to a declaration in the repository source.

The construction follows one pattern: state algebraic content as predicates
over raw matrices, move to subtypes exactly where an API promises a state,
bundle the carrier and the projection so that range and uniqueness theorems
are exact, and hand finished objects to Mathlib interfaces through small
adapters. The pattern is independent of measurement theory and applies to
other finite-dimensional operator-algebra developments over Mathlib.

The companion \texttt{Dynamics/ChoiCPTP.lean} establishes complete positivity
and trace preservation for both expectations under a custom finite-matrix
\texttt{IsCPTP} interface. There is no theorem connecting the supplied
partition to Physlib's spectral pinching and no rank-one classical
specialization. The CHSH saturation witness rests on declared data: no
theorem produces the state or the events from a committed source run. The
ideal static phase-fit model has an explicit synthetic inhabitant and is not
an operational instrument. The entire
development is finite-dimensional and says nothing about normal maps on
infinite-dimensional von Neumann algebras.

The proved surface contains the arbitrary-partition pinching with its exact
range and bimodule law, the averaging expectation with its tower and
classical-conditioning laws, typed L\"uders naturality, the complete
partition-relative quotient, the supplied-state Robertson inequality, the
state-level CHSH corollary, the exact saturation receipt with its
record-diagonal delimitation, and the ideal phase-fit model with its synthetic
non-discharge counterreceipt; the
companion module adds the custom CPTP proofs.
It contains no Physlib
spectral-partition correspondence and no source-selected quantum instrument;
the two-qubit saturation witness consists of declared data, and the
operational phase instrument and its source custody are not constructed.

\section*{Declarations}

\textbf{Funding.} No funding was received for this work.

\textbf{Competing interests.} The author declares no competing interests.

\textbf{Author contributions.} Bernhard Mueller is the sole author, designed
the study, reviewed the formal development, and takes responsibility for the
manuscript and artifact.

\textbf{Data and code availability.} No empirical data were generated. The
canonical Lean source is available in the Observer Patch Holography
repository \cite{observerPatchHolographyRepo}. The repository revision
containing this manuscript provides the Lean source, pinned dependency files,
build instructions, and declaration-level axiom queries under a Creative
Commons BY-NC-SA 4.0 license. No persistent archive identifier is claimed.

\bibliographystyle{plainnat}
\bibliography{machine_checked_finite_event_algebras}

\section*{AI Assistance Disclosure}
This research project used research-grade commercial models, including
Anthropic's Fable and OpenAI's GPT-5.6-Sol, for research support, software
development, editing, and synthesis. The authors are responsible for the
paper's claims, methods, and final text.

\end{document}
